\ifdefined\gkver\else\def\gkver{student}\fi
\documentclass[\gkver]{gaokaozhenti}
\begin{document}

\chapter{2012年全国新课标}
%% paperType: 理综物理部分

\begin{enumerate}
\item
%% number: 14
%% paperName: 2012全国新课标第 14 题 6 分
%% typeId: 2
%% type: 多选题
%% chapter: 运动和力
%% point: 牛顿第一定律
%% method: 无
%% score: 6
%% degree: 650
%% duplicateId: 0
%% body:
伽利略根据小球在斜面上运动的实验和理想实验，提出了惯性的概念，从而奠定了牛顿力学的基础。早期物理学家关于惯性有下列说法，其中正确的是\xzanswer{AD}
\fourchoices[answer=AD]
{物体抵抗运动状态变化的性质是惯性}
{没有力作用，物体只能处于静止状态}
{行星在圆周轨道上保持匀速率运动的性质是惯性}
{运动物体如果没有受到力的作用，将继续以同一速度沿同一直线运动}
%% answer: AD
%% memo:
\memoanswer{
惯性是物体抵抗运动状态变化而保持静止或匀速直线运动状态的性质，A 正确；没有力的作用，物体将处于静止或匀速直线运动状态，B 错误；行星在圆形轨道上保持匀速率运动的原因是行星受到地球的万有引力作用，不是由于惯性，C 错误；运动物体如果没有受到力的作用，将一直匀速直线运动下去，D 正确。
}

\item
%% number: 15
%% paperName: 2012全国新课标第 15 题 6 分
%% typeId: 2
%% type: 多选题
%% chapter: 曲线运动
%% point: 平抛运动
%% method: 无
%% score: 6
%% degree: 650
%% duplicateId: 0
%% body:
如图，$x$ 轴在水平地面内，$y$ 轴沿竖直方向。图中画出了从 $y$ 轴上沿 $x$ 轴正向抛出的三个小球 a、b 和 c 的运动轨迹，其中 b 和 c 是从同一点抛出的，不计空气阻力，则\xzanswer{BD}
\onepicture{figs/15.png}
\fourchoices[answer=BD]
{a 的飞行时间比 b 的长}
{b 和 c 的飞行时间相同}
{a 的水平速度比 b 的小}
{b 的初速度比 c 的大}
%% answer: BD
%% memo:
\memoanswer{
平抛运动可看成水平方向的匀速直线运动和竖直方向的自由落体运动的合运动，因 $y = \frac{1}{2}gt^{2}$，$y_{a} < y_{b} = y_{c}$，所以 b 和 c 飞行时间相等且比 a 的飞行时间长，A 错误，B 正确；因 $x = vt$，$x_{a} > x_{b} > x_{c}$，$t_{a} < t_{b} = t_{c}$，故 $v_{a} > v_{b} > v_{c}$，C 错误，D 正确。
}

\item
%% number: 16
%% paperName: 2012全国新课标第 16 题 6 分
%% typeId: 1
%% type: 单选题
%% chapter: 力物体的平衡
%% point: 动态平衡
%% method: 无
%% score: 6
%% degree: 650
%% duplicateId: 0
%% body:
如图，一小球放置在木板与竖直墙面之间。设墙面对球的压力大小为 $N_{1}$，球对木板的压力大小为 $N_{2}$。以木板与墙连接点所形成的水平直线为轴，将木板从图示位置开始缓慢地转到水平位置。不计摩擦，在此过程中\xzanswer{B}
\onepicture{figs/16.png}
\fourchoices[answer=B]
{$N_{1}$ 始终减小，$N_{2}$ 始终增大}
{$N_{1}$ 始终减小，$N_{2}$ 始终减小}
{$N_{1}$ 先增大后减小，$N_{2}$ 始终减小}
{$N_{1}$ 先增大后减小，$N_{2}$ 先减小后增大}
%% answer: B
%% memo:
\memoanswer{
\onepicture[h=3.4cm, align=right]{figs/16_an.png}

受力分析如图所示：重力的大小方向都不变，可知 $N_{1}$、$N_{2}$ 的合力大小、方向都不变，当木板向下转动时，$N_{1}$、$N_{2}$ 变化如图所示，即 $N_{1}$、$N_{2}$ 都减小，所以正确选项为 B。
}

\item
%% number: 17
%% paperName: 2012全国新课标第 17 题 6 分
%% typeId: 1
%% type: 单选题
%% chapter: 交流电
%% point: 变压器
%% method: 无
%% score: 6
%% degree: 650
%% duplicateId: 0
%% body:
自耦变压器铁芯上只绕有一个线圈，原、副线圈都只取该线圈的某部分，一升压式自耦调压变压器的电路如图所示，其副线圈匝数可调。已知变压器线圈总匝数为 1900 匝；原线圈为 1100 匝，接在有效值为 $220\UV$ 的交流电源上。当变压器输出电压调至最大时，负载 $R$ 上的功率为 $2.0\UkW$。设此时原线圈中电流有效值为 $I_{1}$，负载两端电压的有效值为 $U_{2}$，且变压器是理想的，则 $U_{2}$ 和 $I_{1}$ 分别约为\xzanswer{B}
\onepicture{figs/17.png}
\fourchoices[answer=B]
{$380\UV$ 和 $5.3\UA$}
{$380\UV$ 和 $9.1\UA$}
{$240\UV$ 和 $5.3\UA$}
{$240\UV$ 和 $9.1\UA$}
%% answer: B
%% memo:
\memoanswer{
由理想变压器原理：$\frac{U_{1}}{U_{2}} = \frac{n_{1}}{n_{2}}$ 得：$U_{2} = \frac{n_{2}}{n_{1}}U_{1} = \frac{1900}{1100}\times 220\UV = 380\UV$，又 $P_{1} = P_{2} = U_{1}I_{1} = U_{2}I_{2}$，所以 $I_{1} = \frac{P_{2}}{U_{1}} = \frac{2000}{220}\UA = 9.1\UA$，所以 B 正确。
}

\item
%% number: 18
%% paperName: 2012全国新课标第 18 题 6 分
%% typeId: 2
%% type: 多选题
%% chapter: 电场
%% point: 带电粒子的直线运动
%% method: 无
%% score: 6
%% degree: 650
%% duplicateId: 0
%% body:
如图，平行板电容器的两个极板与水平地面成一角度，两极板与一直流电源相连。若一带电粒子恰能沿图中所示水平直线通过电容器，则在此过程中，该粒子\xzanswer{BD}
\onepicture{figs/18.png}
\fourchoices[answer=BD]
{所受重力与电场力平衡}
{电势能逐渐增加}
{动能逐渐增加}
{做匀变速直线运动}
%% answer: BD
%% memo:
\memoanswer{
\onepicture[h=3.4cm, align=right]{figs/18_an.png}

受力分析如图所示，重力与电场力的合力与速度方向相反，所以粒子做匀减速直线运动，动能减小，所以 A、C 错误，D 正确；因为电场力与速度方向夹角为钝角，所以电场力做负功，电势能增加，即 B 正确。
}

\item
%% number: 19
%% paperName: 2012全国新课标第 19 题 6 分
%% typeId: 1
%% type: 单选题
%% chapter: 电磁感应
%% point: 法拉第电磁感应定律
%% method: 无
%% score: 6
%% degree: 450
%% duplicateId: 0
%% body:
如图，均匀磁场中有一由半圆弧及其直径构成的导线框，半圆直径与磁场边缘重合；磁场方向垂直于半圆面（纸面）向里，磁感应强度大小为 $B_{0}$。使该线框从静止开始绕过圆心 $O$、垂直于半圆面的轴以角速度 $\omega$ 匀速转动半周，在线框中产生感应电流。现使线框保持图中所示位置，磁感应强度大小随时间线性变化。为了产生与线框转动半周过程中同样大小的电流，磁感应强度随时间的变化率 $\frac{\Delta B}{\Delta t}$ 的大小应为\xzanswer{C}
\onepicture{figs/19.png}
\fourchoices[answer=C]
{$\frac{4\omega B_{0}}{\pi}$}
{$\frac{2\omega B_{0}}{\pi}$}
{$\frac{\omega B_{0}}{\pi}$}
{$\frac{\omega B_{0}}{2\pi}$}
%% answer: C
%% memo:
\memoanswer{
线圈匀速转动过程中：$I = \frac{E}{r} = \frac{\frac{1}{2}B_{0}R^{2}\omega}{r} = \frac{1}{2}\frac{B_{0}R^{2}\omega}{r}$；由法拉第电磁感应定律，磁感应强度大小随时间线性变化产生电流 $I = \frac{E}{r} = \frac{1}{r}\frac{\Delta\Phi}{\Delta t} = \frac{1}{r}\frac{\Delta B\cdot\frac{1}{2}\pi R^{2}}{\Delta t} = \frac{1}{2}\pi\frac{1}{r}\frac{\Delta B R^{2}}{\Delta t}$，由题意产生同样大小的电流，所以 $\frac{\Delta B}{\Delta t} = \frac{\omega B_{0}}{\pi}$，故 C 正确。
}

\item
%% number: 20
%% paperName: 2012全国新课标第 20 题 6 分
%% typeId: 1
%% type: 单选题
%% chapter: 电磁感应
%% point: 电磁感应中图象
%% method: 无
%% score: 6
%% degree: 450
%% duplicateId: 0
%% body:
如图，一载流长直导线和一矩形导线框固定在同一平面内，线框在长直导线右侧，且其长边与长直导线平行。已知在 $t=0$ 到 $t=t_{1}$ 的时间间隔内，直导线中电流 $i$ 发生某种变化，而线框中感应电流总是沿顺时针方向；线框受到的安培力的合力先水平向左、后水平向右。设电流 $i$ 正方向与图中箭头方向相同，则 $i$ 随时间 $t$ 变化的图线可能是\xzanswer{A}
\onepicture[align=right]{figs/20.png}
\fourchoices[answer=A, ispicture=true, h=2.2cm]
{figs/20a.png}
{figs/20b.png}
{figs/20c.png}
{figs/20d.png}
%% answer: A
%% memo:
\memoanswer{
由楞次定律可知：线框受力水平向左时，线圈中的磁场要阻碍原磁场引起的磁通量的减弱，说明导线中的电流正在减弱；线框受力水平向右时，线圈中的磁场要阻碍原磁场引起的磁通量的增强，说明导线中的电流正在增强；所以导线中的电流先减弱后增强，所以 C、D 错误；又因线圈中的电流为顺时针方向，所以由右手螺旋定则知线圈产生的磁场垂直纸面向里，因为线圈中的磁场要阻碍原磁场引起的磁通量的减弱，故导线初始状态在导线右侧产生的磁场方向垂直纸面向里，由右手螺旋定则知导线中电流方向为正方向，所以 A 正确，B 错误。
}

\item
%% number: 21
%% paperName: 2012全国新课标第 21 题 6 分
%% typeId: 1
%% type: 单选题
%% chapter: 曲线运动
%% point: 万有引力定律
%% method: 对称法
%% score: 6
%% degree: 250
%% duplicateId: 0
%% body:
假设地球是一半径为 $R$、质量分布均匀的球体。一矿井深度为 $d$。已知质量分布均匀的球壳对壳内物体的引力为零。矿井底部和地面处的重力加速度大小之比为\xzanswer{A}
\fourchoices[answer=A]
{$1-\frac{d}{R}$}
{$1+\frac{d}{R}$}
{$\left(\frac{R-d}{R}\right)^{2}$}
{$\left(\frac{R}{R-d}\right)^{2}$}
%% answer: A
%% memo:
\memoanswer{
在地球表面 $G\frac{Mm}{R^{2}} = mg$，又 $M = \rho\cdot\frac{4}{3}\pi R^{3}$，所以 $g = G\frac{M}{R^{2}} = \frac{4}{3}\pi G\rho R$，因为球壳对球内物体的引力为零，所以在深为 $d$ 的矿井内 $G\frac{Mm}{(R-d)^{2}} = mg'$，得 $g' = G\frac{M}{(R-d)^{2}} = \frac{4}{3}\pi G\rho(R-d)$，所以 $\frac{g'}{g} = \frac{R-d}{R} = 1-\frac{d}{R}$。
}

\item
%% number: 22
%% paperName: 2012全国新课标第 22 题 5 分
%% typeId: 4
%% type: 实验题
%% chapter: 直线运动
%% point: 游标卡尺和螺旋测微仪
%% method: 无
%% score: 5
%% degree: 650
%% duplicateId: 0
%% body:
某同学利用螺旋测微器测量一金属板的厚度。该螺旋测微器校零时的示数如图（a）所示，测量金属板厚度时的示数如图（b）所示。图（a）所示读数为\_\_\_\_\_\_\_\_\_$\Umm$，图（b）所示读数为\_\_\_\_\_\_\_\_\_$\Umm$，所测金属板的厚度为\_\_\_\_\_\_\_\_\_$\Umm$。
\onepicture[align=right]{figs/22.png}
%% answer: 
\jdanswer{
$0.010\Umm$，$6.870\Umm$，$6.860\Umm$
}
%% memo:
\memoanswer{
根据螺旋测微器的读数规则，读数 $=$ 固定刻度 $+$ 可动刻度 $\times 0.01\Umm$，图（a）所示读数为 $0.010\Umm$，图（b）所示读数为 $6.870\Umm$，所测金属板的厚度为 $6.870\Umm - 0.010\Umm = 6.860\Umm$。
}

\item
%% number: 23
%% paperName: 2012全国新课标第 23 题 10 分
%% typeId: 4
%% type: 实验题
%% chapter: 磁场
%% point: 磁感应强度
%% method: 无
%% score: 10
%% degree: 450
%% duplicateId: 0
%% body:
图中虚线框内存在一沿水平方向、且与纸面垂直的匀强磁场。现通过测量通电导线在磁场中所受的安培力，来测量磁场的磁感应强度大小、并判定其方向。所用部分器材已在图中给出，其中 D 为位于纸面内的 U 形金属框，其底边水平，两侧边竖直且等长；$E$ 为直流电源；$R$ 为电阻箱；A 为电流表；S 为开关。此外还有细沙、天平、米尺和若干轻质导线。
\begin{enumerate}
	\item
	在图中画线连接成实验电路图。
	\item
	完成下列主要实验步骤中的填空。

	①按图接线。

	②保持开关 S 断开，在托盘内加入适量细沙，使 D 处于平衡状态；然后用天平称出细沙质量 $m_{1}$。

	③闭合开关 S，调节 $R$ 的值使电流大小适当，在托盘内重新加入适量细沙，使 D\_\_\_\_\_\_\_\_；然后读出\_\_\_\_\_\_\_\_\_\_\_\_\_\_\_\_\_\_，并用天平称出\_\_\_\_\_\_\_\_\_\_\_\_。

	④用米尺测量\_\_\_\_\_\_\_\_\_\_\_\_\_\_\_。
	\item
	用测量的物理量和重力加速度 $g$ 表示磁感应强度的大小，可以得出 $B=$\_\_\_\_\_\_\_\_\_。
	\item
	判定磁感应强度方向的方法是：若\_\_\_\_\_\_\_\_\_\_\_\_，磁感应强度方向垂直纸面向外；反之，磁感应强度方向垂直纸面向里。
\end{enumerate}
\onepicture[align=right]{figs/23.png}
%% answer: 
\jdanswer{
\begin{enumerate}
	\item
	如图所示

	\item
	③重新处于平衡状态，电流表的示数 $I$，此时细沙的质量 $m_{2}$；④D 的底边长度

	\item
	$B=\frac{|m_{2}-m_{1}|g}{Il}$

	\item
	$m_{2}>m_{1}$

\end{enumerate}
}
%% memo:
\memoanswer{
\onepicture[h=3.4cm, align=right]{figs/23_an.png}

（1）依题意：绝缘线圈、电流表、变阻器、开关、电源组成闭合回路，连接如图。

（2）③只有绝缘线圈重新达到平衡状态才可以读出电流表的示数，测出细沙的质量，利用平衡条件得到绝缘线圈所受安培力、拉力、重力的关系。故三处应填：重新处于平衡状态；电流表的示数 $I$；此时细沙的质量。

④要求安培力必须知道磁场中导线受安培力的有效长度，即 D 的底边长度。

（3）由拉力、安培力、重力的平衡得：$F \pm F_{\text{安}} = m_{1}g$，且 $F = m_{2}g$、$F_{\text{安}} = BIl$。所以 $B = \frac{|m_{2}-m_{1}|g}{Il}$。

（4）磁感应强度方向垂直纸面向外时，D 受安培力向下，绝缘线圈平衡必须 $m_{2}>m_{1}$。
}

\item
%% number: 24
%% paperName: 2012全国新课标第 24 题 14 分
%% typeId: 6
%% type: 计算题
%% chapter: 力物体的平衡
%% point: 多力平衡
%% method: 临界
%% score: 14
%% degree: 450
%% duplicateId: 0
%% body:
拖把是由拖杆和拖把头构成的擦地工具（如图）。设拖把头的质量为 $m$，拖杆质量可以忽略；拖把头与地板之间的动摩擦因数为常数 $\mu$，重力加速度为 $g$，某同学用该拖把在水平地板上拖地时，沿拖杆方向推拖把，拖杆与竖直方向的夹角为 $\theta$。
\begin{enumerate}
	\item
	若拖把头在地板上匀速移动，求推拖把的力的大小。
	\item
	设能使该拖把在地板上从静止刚好开始运动的水平推力与此时地板对拖把的正压力的比值为 $\lambda$。已知存在一临界角 $\theta_{0}$，若 $\theta\leq\theta_{0}$，则不管沿拖杆方向的推力多大，都不可能使拖把从静止开始运动。求这一临界角的正切 $\tan\theta_{0}$。
\end{enumerate}
\onepicture[align=right]{figs/24.png}
%% answer: 
\jdanswer{
\begin{enumerate}
	\item
	$\frac{\mu mg}{\sin\theta-\mu\cos\theta}$

	\item
	$\lambda$

\end{enumerate}
}
%% memo:
\memoanswer{
（1）设该同学沿拖杆方向用大小为 $F$ 的力推拖把。将推拖把的力沿竖直和水平方向分解，按平衡条件有

$F\cos\theta+mg=N$ ①

$F\sin\theta=f$ ②

式中 $N$ 和 $f$ 分别为地板对拖把的正压力和摩擦力。按摩擦定律有

$f=\mu N$ ③

联立①②③式得

$F=\frac{\mu}{\sin\theta-\mu\cos\theta}mg$ ④

（2）若不管沿拖杆方向用多大的力都不能使拖把从静止开始运动，应有

$F\sin\theta\leq\lambda N$ ⑤

这时，①式仍满足。联立①⑤式得

$\sin\theta-\lambda\cos\theta\leq\lambda\frac{mg}{F}$ ⑥

现考察使上式成立的 $\theta$ 角的取值范围。注意到上式右边总大于零，且当 $F$ 无限大时极限为零，有

$\sin\theta-\lambda\cos\theta\leq0$ ⑦

使上式成立的 $\theta$ 角满足 $\theta\leq\theta_{0}$，这里 $\theta_{0}$ 是题中所定义的临界角，即当 $\theta\leq\theta_{0}$ 时，不管拖杆方向用多大的力都推不动拖把。临界角的正切为

$\tan\theta_{0}=\lambda$ ⑧
}

\item
%% number: 25
%% paperName: 2012全国新课标第 25 题 18 分
%% typeId: 6
%% type: 计算题
%% chapter: 磁场
%% point: 带电粒子在复合场中的运动
%% method: 无
%% score: 18
%% degree: 450
%% duplicateId: 0
%% body:
如图，一半径为 $R$ 的圆表示一柱形区域的横截面（纸面）。在柱形区域内加一方向垂直于纸面的匀强磁场，一质量为 $m$、电荷量为 $q$ 的粒子沿图中直线在圆上的 a 点射入柱形区域，在圆上的 b 点离开该区域，离开时速度方向与直线垂直。圆心 O 到直线的距离为 $\frac{3}{5}R$。现将磁场换为平行于纸面且垂直于直线的匀强电场，同一粒子以同样速度沿直线在 a 点射入柱形区域，也在 b 点离开该区域。若磁感应强度大小为 $B$，不计重力，求电场强度的大小。
\onepicture[align=right]{figs/25.png}
%% answer: 
\jdanswer{
$\frac{14qRB^{2}}{5m}$
}
%% memo:
\memoanswer{
\onepicture[h=3.0cm, align=right]{figs/25_an.png}

粒子在磁场中做圆周运动。设圆周半径为 $r$，由牛顿第二定律和洛伦兹力公式得

$qvB=m\frac{v^{2}}{r}$ ①

式中 $v$ 为粒子在 a 点的速度。

过 b 点和 O 点作直线的垂线，分别与直线交于 c 和 d 点。由几何关系，线段 $\overline{ab}$、$\overline{bc}$ 和过 a、b 两点的轨迹圆弧的两条半径（未画出）围成一正方形。因此

$\overline{ab}=\overline{bc}=r$ ②

设 $\overline{cd}=x$。由几何关系得

$\overline{ac}=\frac{4}{5}R+x$ ③

$\overline{bc}=\frac{3}{5}R+\sqrt{R^{2}-x^{2}}$ ④

联立②③④式得 $r=\frac{7}{5}R$ ⑤

再考虑粒子在电场中的运动。设电场强度的大小为 $E$，粒子在电场中做类平抛运动。设其加速度大小为 $a$，由牛顿第二定律和带电粒子在电场中的受力公式得

$qE=ma$ ⑥

粒子在电场力方向和直线方向所走的距离均为 $r$，由运动学公式得

$r=\frac{1}{2}at^{2}$ ⑦

$r=vt$ ⑧

式中 $t$ 是粒子在电场中运动的时间。联立①⑤⑥⑦⑧式得

$E=\frac{14}{5}\frac{qRB^{2}}{m}$
}

\item
%% number: 33
%% paperName: 2012全国新课标第 33 题 6 分
%% typeId: 2
%% type: 多选题
%% chapter: 气体性质
%% point: 热力学第一定律
%% method: 无
%% score: 6
%% degree: 650
%% duplicateId: 0
%% body:
关于热力学定律，下列说法正确的是\_\_\_\_\_\_\_\_\_。\xzanswer{ACE}
\fivechoices[answer=ACE]
{为了增加物体的内能，必须对物体做功或向它传递热量}
{对某物体做功，必定会使该物体的内能增加}
{可以从单一热源吸收热量，使之完全变为功}
{不可能使热量从低温物体传向高温物体}
{功转变为热的实际宏观过程是不可逆过程}
%% answer: ACE
%% memo:
\memoanswer{
做功和热传递是改变内能的两种方式，A 正确，B 错误；可以从单一热源吸收热量，使之完全变为功，也可以使热量从低温物体传给高温物体，但要引起其他变化，C 正确，D 错误；根据热力学第二定律，一切与热现象有关的宏观过程都是不可逆的，E 正确。
}

\item
%% number: 33
%% paperName: 2012全国新课标第 33 题 9 分
%% typeId: 6
%% type: 计算题
%% chapter: 气体性质
%% point: 理想气体状态方程
%% method: 无
%% score: 9
%% degree: 450
%% duplicateId: 0
%% body:
如图，由 U 形管和细管连接的玻璃泡 A、B 和 C 浸泡在温度均为 $0\UdC$ 的水槽中，B 的容积是 A 的 3 倍。阀门 S 将 A 和 B 两部分隔开。A 内为真空，B 和 C 内都充有气体。U 形管内左边水银柱比右边的低 $60\Umm$。打开阀门 S，整个系统稳定后，U 形管内左右水银柱高度相等。假设 U 形管和细管中的气体体积远小于玻璃泡的容积。
\begin{enumerate}
	\item
	求玻璃泡 C 中气体的压强（以 $\UmmHg$ 为单位）。
	\item
	将右侧水槽的水从 $0\UdC$ 加热到一定温度时，U 形管内左右水银柱高度差又为 $60\Umm$，求加热后右侧水槽的水温。
\end{enumerate}
\onepicture[align=right]{figs/33.png}
%% answer: 
\jdanswer{
\begin{enumerate}
	\item
	$180\UmmHg$

	\item
	$364\UK$

\end{enumerate}
}
%% memo:
\memoanswer{
（1）在打开阀门 S 前，两水槽水温均为 $T_{0}=273\UK$。设玻璃泡 B 中气体的压强为 $p_{1}$，体积为 $V_{B}$，玻璃泡 C 中气体的压强为 $p_{C}$，依题意有

$p_{1}=p_{C}+\Delta p$ ①

式中 $\Delta p=60\UmmHg$。打开阀门 S 后，两水槽水温仍为 $T_{0}$，设玻璃泡 B 中气体的压强为 $p_{B}$。依题意，有 $p_{B}=p_{C}$ ②

玻璃泡 A 和 B 中气体的体积为

$V_{2}=V_{A}+V_{B}$ ③

根据玻意耳定律得

$p_{1}V_{B}=p_{B}V_{2}$ ④

联立①②③④式，并代入题给数据得

$p_{C}=\frac{V_{B}}{V_{A}}\Delta p=180\UmmHg$ ⑤

（2）当右侧水槽的水温加热至 $T'$ 时，U 形管左右水银柱高度差为 $\Delta p$。玻璃泡 C 中气体的压强为 $p_{C}'=p_{B}+\Delta p$ ⑥

玻璃泡 C 的气体体积不变，根据查理定律得

$\frac{p_{C}}{T_{0}}=\frac{p_{C}'}{T'}$ ⑦

联立②⑤⑥⑦式，并代入题给数据得

$T'=364\UK$ ⑧
}

\item
%% number: 34
%% paperName: 2012全国新课标第 34 题 6 分
%% typeId: 3
%% type: 填空题
%% chapter: 机械振动
%% point: 振动图象
%% method: 无
%% score: 6
%% degree: 650
%% duplicateId: 0
%% body:
一简谐横波沿 $x$ 轴正向传播，$t=0$ 时刻的波形如图（a）所示，$x=0.30\Um$ 处的质点的振动图线如图（b）所示，该质点在 $t=0$ 时刻的运动方向沿 $y$ 轴\_\_\_\_\_\_\_\_\_（填“正向”或“负向”）。已知该波的波长大于 $0.30\Um$，则该波的波长为\_\_\_\_\_\_\_\_\_$\Um$。
\onepicture[align=right]{figs/34.png}
%% answer: 
\jdanswer{
正向，$0.8\Um$
}
%% memo:
\memoanswer{
（1）由 b 图可知，$0$ 时刻质点振动方向沿 $y$ 轴正向；根据质点带动法和波向右传播，由 a 图知介质中各质点的振动方向如图所示，由振动方程 $y=A\sin\frac{2\pi}{T}t$，即 $\sqrt{2}=2\sin\frac{2\pi}{T}t$，得 $\sin\frac{2\pi}{T}t=\frac{\sqrt{2}}{2}$，又波长大于 $0.3\Um$，所以 $\frac{2\pi}{T}t=\frac{3}{4}\pi$，得 $t=\frac{3}{8}T$，又 $v=\frac{\lambda}{T}=\frac{\Delta x}{\Delta t}=\frac{0.3}{\frac{3}{8}T}$，所以 $\lambda=0.8\Um$。
}

\item
%% number: 34
%% paperName: 2012全国新课标第 34 题 9 分
%% typeId: 6
%% type: 计算题
%% chapter: 光学
%% point: 光的折射
%% method: 无
%% score: 9
%% degree: 650
%% duplicateId: 0
%% body:
一玻璃立方体中心有一点状光源。今在立方体的部分表面镀上不透明薄膜，以致从光源发出的光线只经过一次折射不能透出立方体。已知该玻璃的折射率为 $\sqrt{2}$，求镀膜的面积与立方体表面积之比的最小值。
%% answer: 
\jdanswer{
$\frac{\pi}{4}$
}
%% memo:
\memoanswer{
\onepicture[h=3.0cm, align=right]{figs/34_an.png}

如图，考虑从玻璃立方体中心 O 点发出的一条光线，假设它斜射到玻璃立方体上表面发生折射。根据折射定律有

$n\sin\theta=\sin\alpha$ ①

式中，$n$ 是玻璃的折射率，入射角等于 $\theta$，$\alpha$ 是折射角。

现假设 A 点是上表面面积最小的不透明薄膜边缘上的一点。由题意，在 A 点刚好发生全反射，故

$\alpha_{A}=\frac{\pi}{2}$ ②

设线段 OA 在立方体上表面的投影长为 $R_{A}$，由几何关系有

$\sin\theta_{A}=\frac{R_{A}}{\sqrt{R_{A}^{2}+\left(\frac{a}{2}\right)^{2}}}$ ③

式中 $a$ 为玻璃立方体的边长。由①②③式得

$R_{A}=\frac{a}{2\sqrt{n^{2}-1}}$ ④

由题给数据得 $R_{A}=\frac{a}{2}$ ⑤

由题意，上表面所镀的面积最小的不透明薄膜应是半径为 $R_{A}$ 的圆。所求的镀膜面积 $S'$ 与玻璃立方体的表面积 $S$ 之比为

$\frac{S'}{S}=\frac{6\pi R_{A}^{2}}{6a^{2}}$ ⑥

由⑤⑥式得 $\frac{S'}{S}=\frac{\pi}{4}$ ⑦
}

\item
%% number: 35
%% paperName: 2012全国新课标第 35 题 6 分
%% typeId: 3
%% type: 填空题
%% chapter: 原子和原子核
%% point: 质能方程
%% method: 无
%% score: 6
%% degree: 650
%% duplicateId: 0
%% body:
氘核和氚核可发生热核聚变而释放巨大的能量，该反应方程为：$\ce{^{2}_{1}H + ^{3}_{1}H -> ^{4}_{2}He + x}$，式中 $x$ 是某种粒子。已知：$\ce{^{2}_{1}H}$、$\ce{^{3}_{1}H}$、$\ce{^{4}_{2}He}$ 和粒子 $x$ 的质量分别为 $2.0141\Uu$、$3.0161\Uu$、$4.0026\Uu$ 和 $1.0087\Uu$；$1\Uu=931.5\UMeVcSq$，$c$ 是真空中的光速。由上述反应方程和数据可知，粒子 $x$ 是\_\_\_\_\_\_\_\_\_\_，该反应释放出的能量为\_\_\_\_\_\_\_\_\_$\UMeV$（结果保留 3 位有效数字）。
%% answer: 
\jdanswer{
$\ce{^{1}_{0}n}$（或中子），$17.6\UMeV$
}
%% memo:
\memoanswer{
根据 $\ce{^{2}_{1}H + ^{3}_{1}H -> ^{4}_{2}He + x}$ 并结合质量数守恒和电荷数守恒知 $x$ 为 $\ce{^{1}_{0}n}$。由质能方程 $\Delta E=\Delta mc^{2}$ 得：

$\Delta E=(m_{\ce{^{2}_{1}H}}+m_{\ce{^{3}_{1}H}}-m_{\ce{^{4}_{2}He}}-m_{\ce{^{1}_{0}n}})c^{2}=(2.0141+3.0161-4.0026-1.0087)\times 931.5\UMeV=17.6\UMeV$。
}

\item
%% number: 35
%% paperName: 2012全国新课标第 35 题 9 分
%% typeId: 6
%% type: 计算题
%% chapter: 运动和力
%% point: 动量守恒定律
%% method: 无
%% score: 9
%% degree: 450
%% duplicateId: 0
%% body:
如图，小球 a、b 用等长细线悬挂于同一固定点 O。让球 a 静止下垂，将球 b 向右拉起，使细线水平。从静止释放球 b，两球碰后粘在一起向左摆动，此后细线与竖直方向之间的最大偏角为 $60^{\circ}$。忽略空气阻力，求：
\begin{enumerate}
	\item
	两球 a、b 的质量之比；
	\item
	两球在碰撞过程中损失的机械能与球 b 在碰前的最大动能之比。
\end{enumerate}
\onepicture[align=right]{figs/35.png}
%% answer: 
\jdanswer{
\begin{enumerate}
	\item
	$\sqrt{2}-1$

	\item
	$1-\frac{\sqrt{2}}{2}$

\end{enumerate}
}
%% memo:
\memoanswer{
（i）设球 b 的质量为 $m_{2}$，细线长为 $L$，球 b 下落至最低点、但未与球 a 相碰时的速率为 $v$，由机械能守恒定律得

$m_{2}gL=\frac{1}{2}m_{2}v^{2}$ ①

式中 $g$ 是重力加速度的大小。设球 a 的质量为 $m_{1}$；在两球碰后的瞬间，两球共同速度为 $v'$，以向左为正。由动量守恒定律得

$m_{2}v=(m_{1}+m_{2})v'$ ②

设两球共同向左运动到最高处时，细线与竖直方向的夹角为 $\theta$，由机械能守恒定律得

$\frac{1}{2}(m_{1}+m_{2})v'^{2}=(m_{1}+m_{2})gL(1-\cos\theta)$ ③

联立①②③式得

$\frac{m_{1}}{m_{2}}=\frac{1}{\sqrt{1-\cos\theta}}-1$ ④

代入题给数据得

$\frac{m_{1}}{m_{2}}=\sqrt{2}-1$ ⑤

（ii）两球在碰撞过程中的机械能损失是

$Q=m_{2}gL-(m_{1}+m_{2})gL(1-\cos\theta)$ ⑥

联立①⑥式，$Q$ 与碰前球 b 的最大动能 $E_{k}$（$E_{k}=\frac{1}{2}m_{2}v^{2}$）之比为

$\frac{Q}{E_{k}}=1-\frac{m_{1}+m_{2}}{m_{2}}(1-\cos\theta)$ ⑦

联立⑤⑦式，并代入题给数据得

$\frac{Q}{E_{k}}=1-\frac{\sqrt{2}}{2}$ ⑧
}

\end{enumerate}

\end{document}
